\documentclass[11pt,a4paper]{article}

% ---------------------------------------------------------------------------
% Paquetes
% ---------------------------------------------------------------------------
\usepackage[utf8]{inputenc}
\usepackage[T1]{fontenc}
\usepackage[spanish]{babel}
\usepackage{geometry}
\geometry{a4paper, margin=2.2cm}
\usepackage{amsmath}
\usepackage{amssymb}
\usepackage{graphicx}
\usepackage{booktabs}
\usepackage{float}
\usepackage{caption}
\usepackage{enumitem}
\usepackage{xcolor}
\usepackage{hyperref}
\hypersetup{colorlinks=true, linkcolor=blue, citecolor=blue, urlcolor=blue}
\usepackage{titlesec}
\usepackage{fancyhdr}
\usepackage{parskip}
\usepackage{listings}

\lstset{
  language=Python,
  basicstyle=\ttfamily\footnotesize,
  keywordstyle=\color{blue},
  stringstyle=\color{red},
  commentstyle=\color{green!50!black},
  frame=single,
  breaklines=true,
  showstringspaces=false
}

\titlespacing*{\section}{0pt}{10pt}{6pt}
\titlespacing*{\subsection}{0pt}{8pt}{4pt}

\pagestyle{fancy}
\fancyhf{}
\rhead{\small Algoritmos genéticos para termas solares en Perú}
\lhead{\small Trabajo Final de Curso}
\cfoot{\thepage}

\setlength{\parindent}{0pt}


\begin{document}

% 
\begin{titlepage}
    \centering
    \vspace*{1cm}
    {\Large\bfseries Universidad Nacional de San Agustín de Arequipa\\[0.3em]
    Escuela Profesional de Ciencia de la Computación\\[0.5em]}
    {\Large Curso: Inteligencia Artificial\\[2cm]}
    \begin{figure}[H]
        \centering
        \includegraphics[width=0.2\textwidth]{logo_unsa.png} % ESPACIO PARA LOGO
        \label{fig:logo}
    \end{figure}
    \rule{\linewidth}{0.5pt}\\[0.5em]
    {\LARGE\bfseries Réplica basada en la tesis de Ingeniería Mecatrónica\\[0.3em]}
    {\Large Algoritmos Genéticos para Determinar los Ángulos de Inclinación\\
Óptimos de Tubos al Vacío de Termas Solares en las\\
Principales Ciudades del Perú \\[0.5em]}
    {\Large\bfseries Presentado por:\par}
        {\Large Mamani Yucra Edilson Bonet\par}
    {\Large\bfseries Docente:\par}
        {\Large Rolando Cardenas Talavera\par}
    \rule{\linewidth}{0.5pt}\\[3cm]
    {\large Arequipa, \today}
    \vfill
\end{titlepage}

\section{Introducción}
% =============================================================================

El Perú es uno de los países con mayor potencial de radiación solar del planeta,
alcanzando niveles de insolación extremadamente altos en gran parte de su
territorio. Este recurso es aprovechado ampliamente por las termas solares de
tubos al vacío, sistemas de bajo costo utilizados para el calentamiento
doméstico de agua. Sin embargo, la eficiencia de estos sistemas depende
críticamente del ángulo de inclinación con el que se instalan los colectores,
un parámetro que suele fijarse de manera empírica o siguiendo recomendaciones
genéricas del fabricante, sin considerar las particularidades geográficas de
cada ciudad (latitud, altitud y perfil de radiación anual).

La tesis de referencia, desarrollada en el área de Ingeniería Mecatrónica,
aborda este problema aplicando \textbf{Algoritmos Genéticos (AG)} —una técnica
de Inteligencia Artificial bio-inspirada en la selección natural— para
determinar de forma automática el ángulo de inclinación que maximiza la
energía solar anual recolectada por un tubo al vacío, en seis de las
principales ciudades del Perú: Trujillo, Piura, Iquitos, Tacna, Arequipa y
Puno.

El presente trabajo tiene como objetivo replicar íntegramente esta
metodología \textbf{desde cero}, sin el uso de ninguna librería de
Inteligencia Artificial u optimización (ni DEAP, ni \texttt{scipy}, ni
similares), implementando manualmente tanto el modelo matemático de geometría
solar que define la función objetivo, como el propio Algoritmo Genético con
sus operadores de selección, cruzamiento y mutación. Adicionalmente, se
propone e implementa una \textbf{mejora} justificada al algoritmo original,
convirtiéndolo en un algoritmo memético, y se comparan los resultados
obtenidos contra los reportados por el autor de la tesis.

% =============================================================================
\section{Descripción del problema real resuelto}
% =============================================================================

El problema real consiste en \textbf{maximizar la energía solar anual
recolectable por unidad de longitud de un tubo al vacío}, variando su único
parámetro de diseño libre: el ángulo de inclinación $\beta$ (en grados).

Un tubo al vacío recibe radiación solar directa y difusa a lo largo del año.
La cantidad de energía captada depende de una geometría solar compleja que
involucra:

\begin{itemize}[nosep]
  \item La posición del sol en cada instante del día y del año (declinación,
        ángulo horario, ángulo cenital).
  \item La atenuación atmosférica de la radiación (modelo de cielo despejado
        de Hottel).
  \item El sombreado mutuo entre tubos adyacentes, que depende del ángulo de
        inclinación $\beta$, el diámetro de los tubos ($D_1=47$mm interior,
        $D_2=58$mm cobertor) y la distancia de separación entre sus centros
        ($B$).
\end{itemize}

Existe una relación no lineal y no monótona entre $\beta$ y la energía anual
recolectada: ángulos muy bajos maximizan la captación al mediodía pero
desaprovechan la radiación de mañana/tarde; ángulos muy altos reducen el
área efectiva de captación y aumentan el autosombreado. Esto genera una
función objetivo \textbf{unimodal} (un único máximo) en el rango de interés,
razón por la cual el problema es un candidato natural —y a la vez
relativamente sencillo, al tratarse de una sola variable de decisión
continua— para ser resuelto mediante optimización evolutiva.

La función objetivo replicada corresponde a la ecuación (34) de la tesis:
la \textbf{energía anual acumulada} $H_{anual}$, obtenida sumando, para los
365 días del año, la integral diaria de la potencia total (directa + difusa)
incidente sobre el tubo, entre la hora del amanecer y la del atardecer.

Se replicaron los datos geográficos exactos (latitud, longitud y altitud,
tomados en la plaza de armas) de las seis ciudades analizadas en la tesis,
así como la geometría de los tubos ($D_1=47$mm, $D_2=58$mm, $B=80$mm, valor
típico usado por la tesis).

\begin{table}[H]
\centering
\small
\caption{Dataset de coordenadas geográficas y parámetros de los tubos al vacío.}
\begin{tabular}{lccc}
\toprule
\textbf{Ciudad} & \textbf{Latitud (°)} & \textbf{Longitud (°)} & \textbf{Altitud (km)} \\
\midrule
Trujillo & -8.1117 & -79.0286 & 0.034 \\
Tacna & -18.0137 & -70.2510 & 0.552 \\
Arequipa & -16.4154 & -71.5218 & 2.335 \\
Piura & -5.1971 & -80.6266 & 0.055 \\
Puno & -15.8405 & -70.0279 & 3.827 \\
Iquitos & -3.7493 & -73.2444 & 0.106 \\
\bottomrule
\multicolumn{4}{l}{\textbf{Parámetros fijos del tubo al vacío:}} \\
\multicolumn{4}{l}{Diámetro interno ($D_1$) = 47 mm, Diámetro externo ($D_2$) = 58 mm} \\
\multicolumn{4}{l}{Longitud ($L_{tubo}$) = 1.80 m, Separación entre centros ($B$) = 80 mm} \\
\end{tabular}
\end{table}

% =============================================================================
\section{Explicación del algoritmo original (de la tesis)}
% =============================================================================

\subsection{Función objetivo y modelo matemático}

El modelo matemático de geometría solar se implementó desde cero (usando
únicamente el módulo estándar \texttt{math} de Python), reproduciendo con
fidelidad las ecuaciones descritas en la tesis, basadas en Duffie, Beckman \&
Blair (2020) y en Tang, Gao \& Yu (2009):

\begin{align}
\delta &= \text{declinación solar (polinomio de Spencer, 1971)}\\
\omega_{ss} &= \min\!\Big(\omega_s,\ \cos^{-1}\!\big(-\tan\delta\tan(\phi+\beta)\big)\Big)
  \quad \text{(atardecer, plano inclinado, hemisferio sur)}\\
\tau_b &= r_0 a_0 + r_1 a_1 \exp\!\Big(\dfrac{-r_k k}{\cos\theta_Z}\Big)
  \quad \text{(transmitancia directa, modelo de Hottel)}\\
G_{bn} &= G_{on}\cos\theta_Z\,\tau_b \, ,\quad
   \tau_d = 0.271-0.294\tau_b \, ,\quad G_d = G_{on}\cos\theta_Z\,\tau_d\\
P_{total} &= D_1 L_{tubo}\Big(G_{bn}\cos\theta_t\, f(\Omega) + \pi\, G_{d\beta}\, F\Big)
  \quad \text{(ecuación 32 de la tesis)}\\
H_{anual} &= \sum_{n=1}^{365}\int_{-\omega_{ss}}^{\omega_{ss}} P_{total}(t)\, dt
  \quad \text{(ecuación 34 de la tesis, función objetivo)}
\end{align}

donde $f(\Omega)$ es la función de aceptación angular y $F$ el factor de
forma difuso, ambos dependientes de los ángulos críticos $\Omega_0,\Omega_1$
determinados por la geometría de los tubos (ecuaciones 25-31 de la tesis).
La integral diaria se evaluó numéricamente mediante la \textbf{regla de
Simpson compuesta}, implementada también desde cero. A continuación se
muestra la función principal en Python que calcula la energía anual iterando
sobre los 365 días del año:

\begin{lstlisting}
def energia_anual_MJ(phi_deg, altitud_km, beta_deg, B_mm, paso_dias=1, n_pasos=241):
    total = 0.0
    dias = range(1, 366, paso_dias)
    for n in dias:
        total += energia_diaria_MJ(n, phi_deg, altitud_km, beta_deg, B_mm, n_pasos=n_pasos)
    if paso_dias > 1:
        total *= 365.0 / len(dias)
    return total
\end{lstlisting}

\subsection{Representación y operadores del Algoritmo Genético}

Siguiendo la sección 3.4.2 de la tesis (que reporta el uso de la librería
DEAP), se implementó un AG con las siguientes características, replicadas
manualmente en Python puro (módulo \texttt{random} únicamente):

\begin{table}[H]
\centering
\small
\caption{Parámetros configurados para el Algoritmo Genético original.}
\begin{tabular}{ll}
\toprule
\textbf{Parámetro} & \textbf{Valor / Método} \\
\midrule
Representación & Valor real (ángulo $\beta$ en grados) \\
Función de aptitud & Energía anual $H_{anual}(\beta)$ (a maximizar) \\
Tamaño de población & 50 individuos \\
Número de generaciones & 15 \\
Método de selección & Torneo (tamaño $k=3$) \\
Método de cruzamiento & Blend Crossover ($\alpha=0.5$), probabilidad $p_c=0.7$ \\
Método de mutación & Gaussiana ($\sigma=10\%$), probabilidad $p_m=0.1$ \\
Reemplazo & Generacional puro (sin elitismo) \\
\bottomrule
\end{tabular}
\end{table}

El rango de búsqueda de $\beta$ se acotó, para cada ciudad, según la región
donde la propia tesis (Tablas 5 y 6) reportó la mayor energía anual mediante
el barrido inicial de ángulos (0°-35°). El siguiente código muestra cómo se
implementó la selección por torneo y la mutación en el algoritmo puro:

\begin{lstlisting}
def seleccion_torneo(poblacion, k=3):
    seleccionados = []
    for _ in range(len(poblacion)):
        torneo = random.sample(poblacion, k)
        ganador = max(torneo, key=lambda ind: ind["fitness"])
        seleccionados.append(ganador.copy())
    return seleccionados

def mutacion_gaussiana(ind, lo, hi, sigma, prob_mut=0.1):
    if random.random() < prob_mut:
        ind["beta"] += random.gauss(0, sigma)
        # Asegurar que no exceda los limites
        ind["beta"] = max(lo, min(hi, ind["beta"]))
\end{lstlisting}

% =============================================================================
\section{Explicación de la mejora implementada}
% =============================================================================

\subsection{Motivación}

El propio autor de la tesis, en la sección 3.4.4 (``Evaluación de
Resultados''), valida sus resultados del AG comparándolos contra un método
numérico clásico: la \textbf{Búsqueda de la Sección Áurea}
(\emph{Golden Section Search}), un método de optimización unidimensional que
converge de forma muy precisa cuando la función es unimodal en el intervalo
de búsqueda —exactamente el caso de este problema—. Esto sugiere una mejora
natural: en lugar de usar la búsqueda áurea solo como verificación externa
posterior, \textbf{integrarla dentro del propio ciclo evolutivo}, dando
lugar a un \textbf{Algoritmo Memético} (Moscato, 1989): la combinación de un
algoritmo evolutivo (exploración global) con una búsqueda local exacta
(explotación fina).

Adicionalmente, se identificó una debilidad metodológica en el AG original:
al usar reemplazo generacional puro (sin elitismo), el mejor individuo
encontrado en una generación \textbf{puede perderse} si no sobrevive la
selección, el cruce o la mutación de la siguiente generación —un problema
conocido de los AG generacionales simples.

\subsection{Componentes de la mejora propuesta}

\begin{enumerate}[nosep]
  \item \textbf{Elitismo}: se conservan siempre los 2 mejores individuos de
        cada generación, reinsertándolos en el lugar de los 2 peores
        individuos de la descendencia. Esto garantiza una convergencia
        \textbf{monótona no decreciente} del mejor fitness.
  \item \textbf{Mutación gaussiana adaptativa}: la desviación estándar
        $\sigma$ de la mutación decrece linealmente a lo largo de las
        generaciones (de $\sigma_{inicial}=10\%$ del rango a
        $\sigma_{final}=1.5\%$), pasando de una fase de exploración amplia a
        una fase de explotación fina, en lugar de mantener una $\sigma$ fija
        durante las 15 generaciones.
  \item \textbf{Refinamiento memético (Lamarckiano)}: cada 5 generaciones, se
        aplica la Búsqueda de la Sección Áurea en un entorno pequeño
        ($\pm 8\%$ del rango total) alrededor del mejor individuo de la
        población, y si el resultado refinado mejora el fitness, reemplaza
        directamente al individuo en la población (aprendizaje
        Lamarckiano).
\end{enumerate}

Esta combinación es un enfoque estándar y ampliamente documentado en la
literatura de optimización evolutiva para problemas continuos de baja
dimensionalidad, y tiene la ventaja adicional de ser \emph{coherente} con la
propia validación metodológica que el autor de la tesis ya realizó.

% =============================================================================
\section{Resultados obtenidos}
% =============================================================================

Ambos algoritmos (original y mejorado) fueron ejecutados con los mismos
parámetros (población=50, generaciones=15, $p_c=0.7$, $p_m=0.1$) y la misma
semilla aleatoria, para las 6 ciudades de la tesis. El detalle generación
por generación de cada ejecución (fitness del mejor, promedio y peor
individuo por generación) se documenta en el archivo de salida
\texttt{resultados\_algoritmo\_genetico.txt} adjunto al presente informe.

\subsection{Tabla comparativa de ángulos óptimos}

\begin{table}[H]
\centering
\small
\caption{Ángulo óptimo $\beta^*$ (grados) y energía anual (MJ/m) para $B=80$mm,
comparado contra la tesis original.}
\begin{tabular}{lccccc}
\toprule
\textbf{Ciudad} & \textbf{$\beta^*$ Tesis} & \textbf{$\beta^*$ AG Original} & \textbf{$\beta^*$ AG Mejorado} & \textbf{Energía Tesis} & \textbf{Energía (este trabajo)}\\
 & (°) & (°) & (°) & (MJ/m) & (MJ/m)\\
\midrule
Trujillo  & 5.62  & 4.95  & 4.97  & 718.82 & 689.72 \\
Piura     & 3.99  & 3.25  & 3.25  & 725.11 & 695.77 \\
Iquitos   & 2.87  & 2.47  & 2.48  & 730.23 & 700.33 \\
Tacna     & 13.64 & 11.75 & 11.77 & 719.45 & 685.19 \\
Arequipa  & 12.66 & 11.30 & 11.36 & 793.56 & 746.82 \\
Puno      & 12.22 & 11.05 & 10.98 & 804.70 & 755.28 \\
\bottomrule
\end{tabular}
\end{table}

\textbf{Concordancia de ángulos:} la diferencia absoluta promedio entre el
ángulo óptimo hallado por el AG (ambas versiones) y el reportado por la
tesis es de aproximadamente $1.1°$, con la ciudad de Iquitos mostrando la
mayor cercanía ($-0.40°$) y Tacna la mayor discrepancia ($-1.89°$). En las 6
ciudades se conserva el mismo orden relativo de los ángulos óptimos
(Iquitos $<$ Piura $<$ Trujillo $<$ Puno $\approx$ Arequipa $<$ Tacna), lo
cual confirma que el modelo matemático reimplementado captura correctamente
la tendencia física esperada: ciudades de mayor latitud requieren mayor
inclinación.

\textbf{Concordancia de energía:} la energía anual obtenida en este trabajo
es sistemáticamente entre un 94\% y un 96\% de la reportada en la tesis
(razón promedio $\approx 0.95$) para las 6 ciudades. Esta diferencia
proporcional y consistente es esperable, dado que el modelo se reimplementó
desde cero a partir de las ecuaciones del documento, sin acceso al código
fuente original ni a los supuestos de calibración numérica no explicitados
por el autor (p.ej. resolución exacta de la integración numérica, manejo de
casos límite del sombreado). El hecho de que la razón sea consistente entre
las 6 ciudades (no aleatoria) sugiere una diferencia sistemática de
calibración más que un error conceptual del modelo.

\subsection{Comparación de convergencia: AG Original vs. AG Mejorado}

La Figura \ref{fig:convergencia} muestra la evolución del mejor y el peor
individuo de la población, generación a generación, para la ciudad de
Trujillo (resultado representativo de las 6 ciudades).

\begin{figure}[H]
\centering
\includegraphics[width=0.75\textwidth]{conver.png}
\caption{Convergencia del AG Original vs. AG Mejorado (Trujillo). El AG
Mejorado mantiene una población mucho más homogénea (peor individuo mucho
más cercano al mejor) gracias al elitismo y la mutación adaptativa.}
\label{fig:convergencia}
\end{figure}

Ambos algoritmos alcanzan prácticamente el mismo valor de fitness máximo (la
diferencia entre el mejor individuo de ambos en la generación 15 es menor a
$0.001\%$), lo cual es esperable dado que el problema es unidimensional y
unimodal —relativamente sencillo para cualquier AG bien configurado—. La
diferencia relevante aparece en la \textbf{robustez de la convergencia de
toda la población}: en el AG Original, el peor individuo de la generación 15
todavía se encuentra a $3.82$ MJ/m del mejor (fitness peor $=687.10$ vs.
mejor $=690.92$), evidenciando que, sin elitismo, una fracción de la
población puede seguir siendo subóptima incluso al final de la ejecución.
En el AG Mejorado, esa misma brecha se reduce a apenas $0.02$ MJ/m
(peor$=690.90$ vs. mejor$=690.92$), es decir, una reducción de más del
\textbf{99\%} en la dispersión final de la población. Este comportamiento se
repitió de forma consistente en las 6 ciudades evaluadas.

\subsection{Verificación cruzada}

Como verificación adicional (análoga a la que realiza la propia tesis en su
sección 3.4.4), se ejecutó un \textbf{barrido exhaustivo} (fuerza bruta,
paso de $0.05°$) sobre el mismo modelo matemático implementado. En las 6
ciudades, el óptimo hallado por barrido coincidió con el hallado por ambas
versiones del AG con una diferencia menor a $0.05°$ (el propio paso de
discretización del barrido), confirmando que ambos algoritmos genéticos
convergieron correctamente al máximo global de la función objetivo
implementada.

% =============================================================================
\section{Conclusiones}
% =============================================================================

\begin{enumerate}[nosep]
  \item Se logró implementar desde cero, sin ninguna librería de
        Inteligencia Artificial ni de optimización, tanto el modelo
        matemático de geometría solar (ecuaciones 1-34 de la tesis) como un
        Algoritmo Genético que replica fielmente la metodología descrita en
        la sección 3.4.2 (población=50, generaciones=15, cruzamiento=0.7,
        mutación=0.1).
  \item El ángulo óptimo de inclinación hallado para las 6 ciudades
        concuerda de forma cercana con el reportado en la tesis (diferencia
        promedio $\approx 1.1°$), y conserva exactamente el mismo orden
        relativo entre ciudades, validando que el problema de optimización
        (una función unimodal de la energía anual en función de $\beta$) fue
        replicado correctamente en su estructura esencial.
  \item La energía anual absoluta calculada difiere de la reportada por el
        autor en un factor sistemático de aproximadamente 0.95, atribuible a
        diferencias de calibración numérica no explicitadas en el documento
        original (y no a un error conceptual, dado que la razón es
        consistente entre ciudades).
  \item La mejora propuesta —un \textbf{algoritmo memético} que combina
        elitismo, mutación gaussiana adaptativa y refinamiento local con
        Búsqueda de la Sección Áurea— no modificó significativamente el
        valor final del óptimo (por tratarse de un problema unidimensional
        relativamente sencillo), pero sí mejoró sustancialmente la
        \textbf{robustez de la convergencia}, reduciendo en más del 99\% la
        dispersión de fitness entre el mejor y el peor individuo de la
        población final.
  \item Esta mejora es especialmente relevante y generalizable a problemas
        de mayor dimensionalidad (por ejemplo, si se optimizara
        simultáneamente el ángulo $\beta$ y la distancia entre tubos $B$),
        donde el elitismo y el refinamiento local memético previenen la
        pérdida de buenas soluciones y aceleran la convergencia sin
        aumentar el número de generaciones.
  \item El trabajo demuestra la viabilidad de aplicar Algoritmos Genéticos,
        implementados completamente desde cero, a un problema real de
        ingeniería mecatrónica/energías renovables, reforzando el carácter
        interdisciplinario de la Inteligencia Artificial como herramienta de
        optimización en dominios de ingeniería física.
\end{enumerate}

% =============================================================================
\section{Referencias bibliográficas}
% =============================================================================

\begin{itemize}[nosep]
  \item Autor de la tesis (2025). \emph{Algoritmos genéticos para determinar
        los ángulos de inclinación óptimos de tubos al vacío de termas
        solares en las principales ciudades del Perú}. Tesis de Ingeniería
        Mecatrónica.
  \item Beckman, W.A., Duffie, J.A., \& Blair, N. (2020). \emph{Solar
        Engineering of Thermal Processes, Photovoltaics and Wind} (5ta ed.).
        Wiley.
  \item Tang, R., Gao, W., \& Yu, Y. (2009). Optimal tilt-angles of
        all-glass evacuated tube solar collectors. \emph{Energy}, 34,
        1387-1395.
  \item Hottel, H.C. (1976). A simple model for estimating the transmittance
        of direct solar radiation through clear atmospheres. \emph{Solar
        Energy}, 18(2), 129-134.
  \item Liu, B.Y.H., \& Jordan, R.C. (1960). The interrelationship and
        characteristic distribution of direct, diffuse and total solar
        radiation. \emph{Solar Energy}, 4(3), 1-19.
  \item Spencer, J.W. (1971). Fourier series representation of the position
        of the sun. \emph{Search}, 2(5), 172.
  \item Moscato, P. (1989). \emph{On Evolution, Search, Optimization,
        Genetic Algorithms and Martial Arts: Towards Memetic Algorithms}.
        Caltech Concurrent Computation Program, Report 826.
  \item Fortin, F.A., De Rainville, F.M., Gardner, M.A., Parizeau, M., \&
        Gagné, C. (2012). DEAP: Evolutionary algorithms made easy.
        \emph{Journal of Machine Learning Research}, 13(1), 2171-2175.
\end{itemize}

\end{document}